\chapter{2014年福建卷（文）}

\section{单选题}

\begin{problem}
若集合 \(P = \{x \mid 2 \leqslant x < 4\}\)，\(Q = \{x \mid x \geqslant 3\}\)，则 \(P \cap Q\) 等于（\quad）

\choices
    {\(\{x \mid 3 \leqslant x < 4\}\)}
    {\(\{x \mid 3 < x < 4\}\)}
    {\(\{x \mid 2 \leqslant x < 3\}\)}
    {\(\{x \mid 2 \leqslant x \leqslant 3\}\)}

\end{problem}

\begin{problem}
复数 \((3 + 2\i)\) 等于（\quad）

\choices
    {\(-2 - 3\i\)}
    {\(-2 + 3\i\)}
    {\(2 - 3\i\)}
    {\(2 + 3\i\)}

\end{problem}

\begin{problem}
以边长为 1 的正方形的一边所在直线为旋转轴，将该正方形旋转一周所得圆柱的侧面积等于（\quad）

\choices
    {\(2\pi\)}
    {\(\pi\)}
    {2}
    {1}

\end{problem}

\begin{problem}
阅读如图所示的程序框图，运行相应的程序，输出的 \(n\) 的值为（\quad）

\choices
    {1}
    {2}
    {3}
    {4}

\end{problem}

\begin{problem}
命题“\(\forall x \in [0, +\infty), x^2 + x \geqslant 0\)”的否定是（\quad）

\choices
    {\(\forall x \in (-\infty, 0), x^2 + x < 0\)}
    {\(\forall x \in (-\infty, 0), x^3 + x \geqslant 0\)}
    {\(\exists x_0 \in [0, +\infty), x_0^2 + x_0 < 0\)}
    {\(\exists x_0 \in [0, +\infty), x_0^2 + x_0 \geqslant 0\)}

\end{problem}

\begin{problem}
已知直线 \(l\) 过圆 \(x^{2} + (y - 3)^{2} = 4\) 的圆心，且与直线 \(x + y + 1 = 0\) 垂直，则 \(l\) 的方程是（\quad）

\choices
    {\(x + y - 2 = 0\)}
    {\(x - y + 2 = 0\)}
    {\(x + y - 3 = 0\)}
    {\(x - y + 3 = 0\)}

\end{problem}

\begin{problem}
将函数 \( y = \sin x \) 的图象向左平移 \( \frac{\pi}{2} \) 个单位，得到函数 \( y = f(x) \) 的函数图象，则下列说法正确的是（\quad）

\choices
    {\(y = f(x)\) 是奇函数}
    {\(y = f(x)\) 的周期是 \(\pi\)}
    {\(y = f(x)\) 的图象关于直线 \(x = \frac{\pi}{2}\) 对称}
    {\(y = f(x)\) 的图象关于点 \(\left(-\frac{\pi}{2},0\right)\) 对称}

\end{problem}

\begin{problem}
若函数 \( y = \log_a x \)（\(a > 0\)，且 \(a \neq 1\)）的图象如图所示，则下列函数图象正确的是（\quad）

\bitmapfigure[max width=0.34\linewidth,max totalheight=4.2cm]{img/2014/fujian_liberal/q8_fig1.png}

\choices
    {\bitmapinclude[max width=3.0cm,max totalheight=4.1cm]{img/2014/fujian_liberal/q8_choice_a.png}}
    {\bitmapinclude[max width=3.0cm,max totalheight=4.1cm]{img/2014/fujian_liberal/q8_choice_b.png}}
    {\bitmapinclude[max width=3.0cm,max totalheight=4.1cm]{img/2014/fujian_liberal/q8_choice_c.png}}
    {\bitmapinclude[max width=3.0cm,max totalheight=4.1cm]{img/2014/fujian_liberal/q8_choice_d.png}}

\end{problem}

\begin{problem}
要制作一个容积为 \(4 \mathrm{~m}^{3}\)，高为 \(1 \mathrm{~m}\) 的无盖长方体容器，已知该容器的底面造价是每平方米 20 元，侧面造价是每平方米 10 元，则该容器的最低总造价是（\quad）

\choices
    {80 元}
    {120 元}
    {160 元}
    {240 元}

\end{problem}

\begin{problem}
设 \(M\) 为平行四边形 \(ABCD\) 对角线交的点，\(O\) 为平行四边形 \(ABCD\) 所在平面内任意一点，则 \(\overrightarrow{OA} + \overrightarrow{OB} + \overrightarrow{OC} + \overrightarrow{OD}\) 等于（\quad）

\choices
    {\(\overrightarrow{OM}\)}
    {\(2\overrightarrow{OM}\)}
    {\(3\overrightarrow{OM}\)}
    {\(4\overrightarrow{OM}\)}

\end{problem}

\begin{problem}
已知圆 \(C: (x - a)^2 + (y - b)^2 = 1\)，平面区域 \(\Omega: \left\{ \begin{array}{l} x + y - 7 \leqslant 0, \\ x - y + 3 \geqslant 0, \\ y \geqslant 0, \end{array} \right.\) 若圆心 \(C \in \Omega\)，且圆 \(C\) 与 \(x\) 轴相切，则 \(a^2 + b^2\) 的最大值为（\quad）

\choices
    {5}
    {29}
    {37}
    {49}

\end{problem}

\begin{problem}
在平面直角坐标系中，两点 \(P_{1}(x_{1}, y_{1})\)，\(P_{2}(x_{2}, y_{2})\) 间的 “L-距离” 定义为 \(\|P_1P_2\| = |x_1 - x_2| + |y_1 - y_2|\)，则平面内与 \(x\) 轴上两个不同的定点 \(F_{1}, F_{2}\) 的 “L-距离” 之和等于定值（大于 \(\|F_1F_2\|\)）的点的轨迹可以是（\quad）

\choices
    {\bitmapinclude[max width=4.1cm,max totalheight=3.2cm]{img/2014/fujian_liberal/q12_choice_a.png}}
    {\bitmapinclude[max width=4.1cm,max totalheight=3.2cm]{img/2014/fujian_liberal/q12_choice_b.png}}
    {\bitmapinclude[max width=4.1cm,max totalheight=3.2cm]{img/2014/fujian_liberal/q12_choice_c.png}}
    {\bitmapinclude[max width=4.1cm,max totalheight=3.2cm]{img/2014/fujian_liberal/q12_choice_d.png}}

\end{problem}

\section{解答题}

\begin{problem}
已知函数 \( f(x) = 2\cos x (\sin x + \cos x) \). (1) 求 \( f\left(\frac{5\pi}{4}\right) \) 的值; (2) 求函数 \( f(x) \) 的最小正周期及单调递增区间.

\end{problem}

\begin{problem}
根据世行 2013 年新标准，人均 GDP 低于 1035 美元为低收入国家; 人均 GDP 为 1035—4085 美元为中等偏下收入国家; 人均 GDP 为 4085—12616 美元为中等偏上收入国家; 人均 GDP 不低于 12616 美元为高收入国家. 某城市有 5 个行政区，各区人口占该城市人口比例及人均 GDP 如下表: 行政区 区人口占城市人口比例 区人均GDP（单位: 美元）A 25\% 8000 B 30\% 4000 C 15\% 6000 D 10\% 3000 E 20\% 10000. (1) 判断该城市人均 GDP 是否达到中等偏上收入国家标准; (2) 现从该城市 5 个行政区中随机抽取 2 个，求抽到的 2 个行政区人均 GDP 都达到中等偏上收入国家标准的概率.

\end{problem}

\begin{problem}
已知曲线 \(F\) 上的点到点 \(F(0,1)\) 的距离比它到直线 \(y = -3\) 的距离小 2. (1) 求曲线 \(F\) 的方程; (2) 曲线 \(F\) 在点 \(P\) 处的切线 \(l\) 与 \(x\) 轴交于点 \(A\). 直线 \(y = 3\) 分别与直线 \(l\) 及 \(y\) 轴交于点 \(M, N\). 以 \(MN\) 为直径作圆 \(C\)，过点 \(A\) 作圆 \(C\) 的切线，切点为 \(B\). 试探究: 当点 \(P\) 在曲线 \(F\) 上运动（点 \(P\) 与原点不重合）时，线段 \(AB\) 的长度是否发生变化\(\perp\) 证明你的结论.

\end{problem}

\begin{problem}
已知函数 \( f(x) = \e^{x} - ax \)（\(a\) 为常数）的图象与 \(y\) 轴交于点 \(A\)，曲线 \( y = f(x) \) 在点 \(A\) 处的切线斜率为 \(-1\). (1) 求 \(a\) 的值及函数 \(f(x)\) 的极值; (2) 证明：当 \(x > 0\) 时，\(x^2 < \e^x\); (3) 证明：对任意给定的正数 \(c\)，总存在 \(x_0\)，使得当 \(x \in (x_0, +\infty)\) 时，恒有 \(x^2 < cx^c\).
\end{problem}

